\documentclass[11pt,a4paper]{ctexart}
% ============================================================
%  函数极限计算与泰勒公式 · 专项练习册
%  编译：xelatex 两遍
%  依据：上海交大高数 A 类期中真题 2008--2025 共 17 套
%        全部极限值、级数系数、条件题导数值均已用 sympy 复核
%        （Homework/_review/tools/verify_practice_limits_taylor.py，36 项全 PASS）
% ============================================================
\usepackage[a4paper,top=2.3cm,bottom=2.3cm,left=2.2cm,right=2.2cm]{geometry}
\usepackage{amsmath,amssymb,amsthm,mathtools,bm}
\usepackage{booktabs,array,longtable,tabularx,multicol}
\usepackage{enumitem}
\usepackage[most]{tcolorbox}
\usepackage{fancyhdr}
\usepackage{hyperref}
\hypersetup{hidelinks,colorlinks=true,linkcolor=blue!50!black,urlcolor=blue!50!black}
\usepackage{titlesec}
\usepackage{pifont}

\linespread{1.28}
\setlength{\parskip}{0.25em}
\setlist[enumerate]{itemsep=4pt,topsep=4pt,parsep=0pt,leftmargin=2.2em}
\titleformat{\section}{\Large\bfseries\color{blue!35!black}}{\thesection}{0.6em}{}
\titleformat{\subsection}{\large\bfseries\color{blue!25!black}}{\thesubsection}{0.5em}{}
\titlespacing*{\section}{0pt}{1.7ex plus .3ex}{1.1ex}
\titlespacing*{\subsection}{0pt}{1.3ex plus .2ex}{0.8ex}
\setcounter{tocdepth}{1}

\pagestyle{fancy}
\fancyhf{}
\fancyhead[L]{\small 函数极限计算与泰勒公式 · 专项练习册}
\fancyhead[R]{\small 陶理 \quad 2026-11-03}
\fancyfoot[C]{\small\thepage}
\renewcommand{\headrulewidth}{0.4pt}

\newcommand{\R}{\mathbb{R}}
\newcommand{\eps}{\varepsilon}
\newcommand{\abs}[1]{\left\lvert #1\right\rvert}
\newcommand{\ee}{\mathrm{e}}
\newcommand{\D}{\mathrm{d}}
\renewcommand{\proofname}{证明}

\newtcolorbox{keypoint}[1][]{enhanced,breakable,colback=blue!3,colframe=blue!40!black,
  fonttitle=\bfseries,title={#1},boxrule=0.5pt,arc=1.5mm,
  left=3mm,right=3mm,top=2mm,bottom=2mm,before skip=4pt,after skip=4pt}
\newtcolorbox{warnbox}[1][]{enhanced,breakable,colback=red!3,colframe=red!50!black,
  fonttitle=\bfseries,title={#1},boxrule=0.5pt,arc=1.5mm,
  left=3mm,right=3mm,top=2mm,bottom=2mm,before skip=4pt,after skip=4pt}
\newtcolorbox{methodbox}[1][]{enhanced,breakable,colback=green!3!white,colframe=green!35!black,
  fonttitle=\bfseries,title={#1},boxrule=0.5pt,arc=1.5mm,
  left=3mm,right=3mm,top=2mm,bottom=2mm,before skip=4pt,after skip=4pt}
\newtcolorbox{notebox}[1][]{enhanced,breakable,colback=gray!5,colframe=gray!55,
  fonttitle=\bfseries,title={#1},boxrule=0.4pt,arc=1mm,
  left=3mm,right=3mm,top=1.5mm,bottom=1.5mm,before skip=3pt,after skip=3pt}

\newenvironment{solution}{\par\smallskip\noindent\textbf{解：}\ignorespaces}{\par\smallskip}
\newcommand{\kw}[1]{\ifmmode\textcolor{blue!45!black}{#1}\else\textbf{\textcolor{blue!45!black}{#1}}\fi}
\newcommand{\yes}{\textcolor{green!50!black}{\ding{51}}}
\newcommand{\no}{\textcolor{red!70!black}{\ding{55}}}
\newcommand{\ans}[1]{\ifmmode\textcolor{red!60!black}{#1}\else\textbf{\textcolor{red!60!black}{#1}}\fi}
\newcommand{\blank}{\underline{\hspace{2.4em}}}

\begin{document}

% ============================================================
\begin{titlepage}
\centering
\vspace*{1.8cm}
{\Huge\bfseries 函数极限计算与泰勒公式\par}
\vspace{0.45cm}
{\Huge\bfseries 专项练习册\par}
\vspace{0.9cm}
{\large A 组 12 题（计算技巧）\quad$\cdot$\quad B 组 13 题（泰勒）\quad$\cdot$\quad 30 分钟限时卷 10 题\par}
\vspace{0.5cm}
{\large 附：泰勒决策卡（一页）\quad$\cdot$\quad 全部答案与解析\par}
\vspace{1.2cm}
{\large 依据 \textbf{上海交通大学 2008--2025 级第一学期高数 A 类期中真题 17 套}\par}
\vspace{0.3cm}
{\large 每题答案均经 sympy 精确计算复核（\texttt{verify\_practice\_limits\_taylor.py}：36 项全通过）\par}
\vfill
{\small\begin{minipage}{0.9\textwidth}
\textbf{今天（11/3，D-7）怎么用}\par\vspace{2pt}
第 1 节是\textbf{泰勒决策卡}，先背下来（10 分钟）；A 组 12 题限时 30 分钟；B 组 13 题限时 40 分钟；
第 4 节限时卷 30 分钟 20 分是今天的\textbf{验收线}。做题顺序：\textbf{先遮住答案独立做，再对解析}。
每道错题写一句「错因 + 触发条件」进错题本 —— 错题本比做完题本身更重要。\par\vspace{3pt}
\textbf{与已有资料的分工}：决策流程与间断点分类见《讲义\_函数极限与间断点》；
第 11 题定义证明见《讲义\_极限定义证明》；本册\textbf{只做题}，不重复讲义里的结论。\par\vspace{3pt}
\textbf{课外补充阅读}（可选，30 分钟选读）：MIT OCW 18.001 第 2 章（Limits、Continuous Functions）与第 10 章（Infinite Series），
见 \texttt{Resources/Downloads/}。
\end{minipage}}
\end{titlepage}

\tableofcontents
\newpage

% ============================================================
\section{泰勒决策卡（一页，闭卷默写）}

\begin{keypoint}[一句话总纲]
\textbf{先问两个问题：(1) 极限是不是 $0/0$（或 $\infty/\infty$）？(2) 分母（被除的那一项）最高是几阶？}
答案是 $k$ 阶，就展到 $k$ 阶；被除的一项是常数，就把分母展到比它高一阶。
\end{keypoint}

\subsection{什么时候用等价无穷小，什么时候必须泰勒}

\begin{center}
\footnotesize
\renewcommand{\arraystretch}{1.3}
\begin{tabularx}{\textwidth}{@{}>{\raggedright\arraybackslash}p{3.9cm} >{\raggedright\arraybackslash}p{5.2cm} X@{}}
\toprule
\textbf{情形} & \textbf{用什么} & \textbf{例子}\\
\midrule
乘积 / 商里的\textbf{因子} & 等价无穷小（$\sin u\sim u$、$\ln(1+u)\sim u$、$1-\cos u\sim\frac{u^2}2$、$(1+u)^\alpha\sim1+\alpha u$） &
$\lim\limits_{x\to0}\dfrac{\sin 3x}{x}=3$\\
\textbf{和差}里要消掉低阶项 & 必须泰勒（或先合并 / 有理化） &
$\lim\limits_{x\to0}\dfrac{\sin x-x}{x^3}=-\frac16$\\
分母是 $x^k$、分子含 $k$ 阶项 & 分子展开到 $k$ 阶，分母展开到 $k-1$ 阶（分母阶数 $=k-1$） &
$\lim\limits_{x\to0}\dfrac{\ee^x-1-x}{x^2}=\frac12$\\
条件给了一个极限，求某阶导数 & 用「$f(x)=f(0)+f'(0)x+\frac{f''(0)}{2}x^2+o(x^2)$」\textbf{比对系数} &
$\lim\limits_{x\to0}\dfrac{f(x)-x}{x^2}=3\Rightarrow f''(0)=6$\\
$1^\infty$ 型 & 取对数化为 $0/0$ 或 $\infty/\infty$ &
$\lim\limits_{x\to0}\left(\dfrac{\sin x}{x}\right)^{1/x^2}=\ee^{-1/6}$\\
\bottomrule
\end{tabularx}
\end{center}

\begin{warnbox}[四个高频坑（考场上最容易丢分的地方）]
\begin{enumerate}[label=(\alph*)]
  \item \textbf{加减式不能直接等价替换}。$\sin x - x$ 里 $\sin x\sim x$，替换后得到 $0$，答案错。
  正确做法：Taylor 到第一个非零阶。
  \item \textbf{分母阶数看错}。分母是 $x^2$ 就要展到 2 阶，不是 1 阶；分母是 $x^3$ 展到 3 阶。
  展多了不会错（只是多写一行），展少了必错。
  \item \textbf{条件题要写成泰勒形式再比系数}。题目给 $\lim\frac{f(x)-g(x)}{x^2}=L$，正确写法是
  $f(x)=g(x)+Lx^2+o(x^2)$，然后逐项比对；不要直接"求导两次"硬套。
  \item \textbf{符号别记错}：$\sqrt{1-x}=1-\frac{x}{2}-\frac{x^2}{8}+\cdots$（负号在 $x^2$ 项），
  $\sqrt{1+x}=1+\frac{x}{2}-\frac{x^2}{8}+\cdots$。这正是 \texttt{7.tex} 第 5(7) 题当初出错的地方。
\end{enumerate}
\end{warnbox}

\begin{notebox}[等价无穷小 3 规则（背下来，30 秒）]
第 1 条：只在\textbf{乘除因子}中替换；第 2 条：加减式先合并、通分、有理化，或直接 Taylor 到第一个非零阶；
第 3 条：做完代 $x=10^{-3}$ 数值抽查符号与量级。
\end{notebox}

% ============================================================
\section{A 组：函数极限计算（12 题，限时 30 分钟）}

\begin{notebox}[A 组的目标]
练\textbf{工具选择}：有理化、提 $\ee$、等价无穷小、夹逼、洛必达。每题先写"选型"，再动笔。
\end{notebox}

\begin{enumerate}[label=\textbf{A\arabic*.}]
  \item $\displaystyle \lim_{x\to0}\dfrac{\sin 3x}{x}$
  \item $\displaystyle \lim_{x\to0}\dfrac{\arctan x}{x}$
  \item $\displaystyle \lim_{x\to0}\dfrac{1-\cos x}{x^{2}}$
  \item $\displaystyle \lim_{x\to0}\dfrac{\ln(1+x)}{x}$
  \item $\displaystyle \lim_{x\to+\infty}\left(1+\dfrac1x\right)^{x}$
  \item $\displaystyle \lim_{x\to0}\dfrac{\sqrt[3]{1+x}-1}{x}$
  \item $\displaystyle \lim_{x\to0}\dfrac{2^{x}-3^{x}}{x}$
  \item $\displaystyle \lim_{x\to0}x\sin\dfrac1x$
  \item $\displaystyle \lim_{x\to0}\dfrac{\sqrt{1+x}-\sqrt{1-x}}{x}$
  \item $\displaystyle \lim_{x\to0}\dfrac{\tan x-\sin x}{x^{3}}$
  \item $\displaystyle \lim_{x\to+\infty}\left(\dfrac{x+1}{x-1}\right)^{x}$
  \item $\displaystyle \lim_{n\to+\infty}\dfrac{1+\frac12+\cdots+\frac1n}{\ln n}$
\end{enumerate}

\vspace{2em}
\noindent\rule{\textwidth}{0.4pt}
\newpage

% ============================================================
\section{B 组：泰勒公式（13 题，限时 40 分钟）}

\begin{notebox}[B 组的目标]
练\textbf{展几阶、展谁、怎么比系数}。每题动笔前先写一行「我要展到 $k$ 阶，因为分母是 $x^k$」或
「题目要 $f^{(n)}(0)$，所以要展到 $n$ 阶」—— 这一行就是步骤分。
\end{notebox}

\begin{enumerate}[label=\textbf{B\arabic*.}]
  \item 写出 $\ln(1+x)$ 的 4 阶麦克劳林展开，并求
        $\displaystyle \lim_{x\to0}\dfrac{\ln(1+x)-x+\frac{x^{2}}2}{x^{3}}$
  \item 求 $\displaystyle \lim_{x\to0}\dfrac{\sin x-x+\frac{x^{3}}6}{x^{5}}$
  \item 求 $\displaystyle \lim_{x\to0}\dfrac{\cos x-1+\frac{x^{2}}2}{x^{4}}$
  \item 求 $\displaystyle \lim_{x\to0}\dfrac{\sqrt{1+x}-1-\frac{x}{2}}{x^{2}}$
  \item 对任意实数 $\alpha$（$\alpha\neq0,1$），求
        $\displaystyle \lim_{x\to0}\dfrac{(1+x)^{\alpha}-1-\alpha x}{\alpha(\alpha-1)x^{2}}$
  \item 设 $f(x)=(1+x)^{\alpha}$，求 $f^{(13)}(0)$ 的表达式
        （提示：$2023$--$2024$ 学年第 15 题考的就是这个）
  \item 求 $f(x)=\ee^{\sin x}$ 的 2 阶麦克劳林展开
  \item 求 $f(x)=\sqrt{x}$ 在 $x=1$ 处的 2 阶泰勒展开
  \item 求 $\displaystyle \lim_{x\to0}\dfrac{\ln(\cos x)}{x^{2}}$
  \item 设 $f(x)$ 在 $x=0$ 处二阶可导，且
        $\displaystyle \lim_{x\to0}\dfrac{f(x)-\ln(1+x)}{x^{2}}=2$，求 $f''(0)$
  \item 求 $\displaystyle \lim_{x\to0}\left(\dfrac{\sin x}{x}\right)^{1/x^{2}}$
  \item 求 $\displaystyle \lim_{x\to0}\dfrac{\ee^{x}-1-x-\frac{x^{2}}2}{x^{3}}$
  \item 设 $f(x)=ax+bx^{2}+o(x^{2})$，且
        $\displaystyle \lim_{x\to0}\dfrac{\ee^{f(x)}-1-x}{x^{2}}=3$，求 $a,b$
\end{enumerate}

\vspace{2em}
\noindent\rule{\textwidth}{0.4pt}
\newpage

% ============================================================
\section{30 分钟限时卷（10 题，20 分）}

\begin{keypoint}[验收线]
\textbf{30 分钟 20 分，低于 14 分}（70\%）就说明"会讲但不会做"，明天 11/4 的主块末尾要加练；
\textbf{低于 10 分}（50\%）说明泰勒还没立住，今晚必须补 §1 决策卡 + B 组错题。
\end{keypoint}

\begin{enumerate}[label=\textbf{L\arabic*.}]
  \item （2 分）$\displaystyle \lim_{x\to0}\dfrac{\tan x-x}{x^{3}}$
  \item （2 分）$\displaystyle \lim_{x\to0^{+}}x^{\sin x}$
  \item （2 分）$\displaystyle \lim_{x\to0}\dfrac{1-\sqrt[3]{1+x}}{x}$
  \item （2 分）$\displaystyle \lim_{x\to+\infty}\dfrac{3x^{2}+2x}{x^{3}-x}$
  \item （2 分）$\displaystyle \lim_{x\to0}\dfrac{\ee^{2x}-\ee^{-2x}}{\sin 2x}$
  \item （2 分）$\displaystyle \lim_{x\to0}\dfrac{\ln(1+2x)}{x}$
  \item （2 分）$\displaystyle \lim_{x\to0}\dfrac{1-\cos 2x}{x\sin x}$
  \item （2 分）$\displaystyle \lim_{x\to0}\left(\dfrac{1+x}{1-x}\right)^{1/x}$
  \item （2 分）$\displaystyle \lim_{x\to0}\dfrac{\cos x-\ee^{-x^{2}/2}}{x^{4}}$
  \item （2 分）设 $f(x)$ 在 $x=0$ 处二阶可导，且
        $\displaystyle \lim_{x\to0}\dfrac{f(x)-f(0)-x}{x^{2}}=3$，求 $f''(0)$
\end{enumerate}

\vspace{2em}
\noindent\rule{\textwidth}{0.4pt}
\newpage

% ============================================================
\section{答案与解析}

\subsection{A 组解析}

\begin{enumerate}[label=\textbf{A\arabic*.},start=1]
  \item 令 $u=3x$，用 $\sin u\sim u$：$\lim\frac{\sin 3x}{x}=3\lim\frac{\sin 3x}{3x}=3$。\ans{$3$}
  \item $\arctan x\sim x$（等价无穷小），故 $\ans{1}$。
  \item $1-\cos x=2\sin^{2}\frac{x}{2}\sim 2\cdot\frac{x^{2}}4=\frac{x^{2}}2$，故 $\ans{\frac12}$。
  \item $\ln(1+x)\sim x$，故 $\ans{1}$。
  \item $1^\infty$ 型，\textbf{取对数}：
        $\ln(\cdot)=x\ln\left(1+\frac1x\right)=x\cdot\frac{\frac1x}{1+o(1)}\to1$，故 $\ans{\ee}$。
  \item 立方差有理化：$\frac{(1+x)-1}{x\left((1+x)^{2/3}+(1+x)^{1/3}+1\right)}\to\frac13$。\ans{$\frac13$}
  \item $\frac{2^{x}-1}{x}-\frac{3^{x}-1}{x}\to\ln2-\ln3=\ln\frac23$。\ans{$\ln\frac23$}
  \item \textbf{有界 $\times$ 无穷小 $=$ 无穷小}：$\abs{\sin\frac1x}\le1$ 而 $x\to0$。\ans{$0$}
  \item 有理化：$\frac{2x}{x\left(\sqrt{1+x}+\sqrt{1-x}\right)}\to\frac22=1$。\ans{$1$}
  \item $\tan x-\sin x=\sin x\cdot\frac{1-\cos x}{\cos x}\sim x\cdot\frac{x^{2}}2\cdot1=\frac{x^{3}}2$。\ans{$\frac12$}
  \item 取对数：$x\left[\ln\left(1+\frac1x\right)-\ln\left(1-\frac1x\right)\right]\to 1+1=2$，故 $\ans{\ee^{2}}$。
        \textbf{提醒}：$2016$--$2017$ 卷考过它的数列版 $\lim\frac{n-\sin n}{n+\sin n}=1$，\textbf{数列不能直接洛必达}。
  \item 夹逼：$\ln(n+1)\le H_n\le1+\ln n$，两边同除 $\ln n\to+\infty$（Stolz 也可）。\ans{$1$}
\end{enumerate}

\subsection{B 组解析}

\begin{enumerate}[label=\textbf{B\arabic*.},start=1]
  \item $\ln(1+x)=x-\frac{x^{2}}2+\frac{x^{3}}3-\frac{x^{4}}4+o(x^{4})$，故
        $\ln(1+x)-x+\frac{x^{2}}2=\frac{x^{3}}3+o(x^{3})$，极限 $\ans{\frac13}$。
        \textbf{注意符号}：$x^{3}$ 项是 $+\frac13$ 不是 $-\frac13$（与 $\ln(1-x)$ 相反）。
  \item $\sin x=x-\frac{x^{3}}6+\frac{x^{5}}{120}+o(x^{5})$，极限 $\ans{\frac1{120}}$。
  \item $\cos x=1-\frac{x^{2}}2+\frac{x^{4}}{24}+o(x^{4})$，极限 $\ans{\frac1{24}}$。
  \item $\sqrt{1+x}=1+\frac{x}{2}-\frac{x^{2}}8+o(x^{2})$，极限 $\ans{-\frac18}$。
        \textbf{这就是 \texttt{7.tex} 第 5(7) 题当年错的符号}。
  \item $(1+x)^{\alpha}=1+\alpha x+\frac{\alpha(\alpha-1)}{2}x^{2}+o(x^{2})$，分子 $=\frac{\alpha(\alpha-1)}2x^{2}+o(x^{2})$，
        除以 $\alpha(\alpha-1)x^{2}$ 得 $\ans{\frac12}$（与 $\alpha$ 无关，只要 $\alpha\neq0,1$）。
  \item 通项 $f^{(n)}(0)=\alpha(\alpha-1)(\alpha-2)\cdots(\alpha-n+1)$，所以
        $\ans{f^{(13)}(0)=\alpha(\alpha-1)(\alpha-2)\cdots(\alpha-12)}$（13 个因子）。
        \textbf{步骤分写法}：先写通项，再代入 $n=13$。
  \item $\sin x=x+o(x^{2})$，$\ee^{u}=1+u+\frac{u^{2}}2+o(u^{2})$，故
        $\ee^{\sin x}=1+\sin x+\frac{\sin^{2}x}2+o(x^{2})=1+x+\frac{x^{2}}2+o(x^{2})$。\ans{$1+x+\frac{x^{2}}2$}
  \item 令 $u=x-1$，$\sqrt{1+u}=1+\frac{u}{2}-\frac{u^{2}}8+o(u^{2})$，故
        $\ans{1+\frac{x-1}{2}-\frac{(x-1)^{2}}8+o((x-1)^{2})}$。
  \item $\cos x=1-\frac{x^{2}}2+o(x^{4})$，$\ln(1+u)=u-\frac{u^{2}}2+o(u^{2})$（$u=-\frac{x^{2}}2$），
        得 $\ln(\cos x)=-\frac{x^{2}}2+o(x^{2})$，极限 $\ans{-\frac12}$。
  \item 条件即 $f(x)=\ln(1+x)+2x^{2}+o(x^{2})=\left(x-\frac{x^{2}}2\right)+2x^{2}+o(x^{2})
        =x+\frac32x^{2}+o(x^{2})$。
        与 $f(x)=f(0)+f'(0)x+\frac{f''(0)}2x^{2}+o(x^{2})$ 比系数：$\frac{f''(0)}2=\frac32$，故 $\ans{f''(0)=3}$。
        \textbf{坑}：$(\ln(1+x))''(0)=-1$，忘了它就会答成 $4$。
  \item $1^\infty$ 型，取对数：$\frac{1}{x^{2}}\ln\frac{\sin x}{x}$，
        而 $\frac{\sin x}{x}=1-\frac{x^{2}}6+o(x^{2})$，故 $\ln(\cdot)=-\frac{x^{2}}6+o(x^{2})$，
        除以 $x^{2}$ 得 $-\frac16$，原式 $\ans{\ee^{-1/6}}$。
  \item $\ee^{x}=1+x+\frac{x^{2}}2+\frac{x^{3}}6+o(x^{3})$，极限 $\ans{\frac16}$。
  \item $f=ax+bx^{2}+o(x^{2})$，$\ee^{f}=1+f+\frac{f^{2}}2+o(x^{3})
        =1+ax+\left(b+\frac{a^{2}}2\right)x^{2}+o(x^{2})$。
        条件 $\frac{\ee^{f}-1-x}{x^{2}}\to3$：$x$ 的系数 $a-1=0\Rightarrow a=1$；
        $x^{2}$ 的系数 $b+\frac{a^{2}}2=3\Rightarrow b=\frac52$。\ans{$a=1,\ b=\frac52$}
\end{enumerate}

\subsection{限时卷解析}

\begin{enumerate}[label=\textbf{L\arabic*.},start=1]
  \item $\tan x=x+\frac{x^{3}}3+o(x^{3})$，极限 $\ans{\frac13}$。
  \item 底数 $\to0$，指数 $\to0$，属 $0^0$ 型：取对数 $\sin x\cdot\ln x\to0\cdot(-\infty)=0$（有界 $\times$ 无穷小），故 $\ans{1}$。
        \textbf{只定义在 $x>0$}，所以是单侧题。
  \item 立方差有理化：$\frac{(1+x)-1}{x\left((1+x)^{2/3}+(1+x)^{1/3}+1\right)}\to-\frac13$。\ans{$-\frac13$}
  \item $\frac{3x^{2}+2x}{x^{3}-x}=\frac{3+\frac2x}{x^{2}-1}\to0$。\ans{$0$}
  \item 分子 $\ee^{2x}-\ee^{-2x}=2\sinh 2x\sim4x$，分母 $\sin2x\sim2x$，故 $\ans{2}$。
  \item $\ln(1+2x)\sim2x$，故 $\ans{2}$。
  \item $1-\cos2x\sim\frac{(2x)^{2}}2=2x^{2}$，$x\sin x\sim x^{2}$，故 $\ans{2}$。
  \item 取对数：$\frac{1}{x}\left[\ln(1+x)-\ln(1-x)\right]\to\frac1x(2x)=2$，故 $\ans{\ee^{2}}$。
  \item 两边都展到 4 阶：$\cos x=1-\frac{x^{2}}2+\frac{x^{4}}{24}+o(x^{4})$，
        $\ee^{-x^{2}/2}=1-\frac{x^{2}}2+\frac{x^{4}}8+o(x^{4})$，
        差 $=\left(\frac1{24}-\frac18\right)x^{4}=-\frac{x^{4}}{12}+o(x^{4})$，极限 $\ans{-\frac1{12}}$。
  \item 条件即 $f(x)=f(0)+x+3x^{2}+o(x^{2})$，与泰勒式比系数：$\frac{f''(0)}2=3$，故 $\ans{f''(0)=6}$。
\end{enumerate}

% ============================================================
\appendix
\section{附录：备用表}

\begin{center}
\footnotesize
\renewcommand{\arraystretch}{1.25}
\begin{tabularx}{\textwidth}{@{}>{\raggedright\arraybackslash}p{4.4cm} X@{}}
\toprule
\textbf{常用展开} & \textbf{公式}\\
\midrule
$\ee^x$ & $1+x+\frac{x^{2}}2+\frac{x^{3}}6+o(x^{3})$\\
$\sin x$ & $x-\frac{x^{3}}6+o(x^{3})$\\
$\cos x$ & $1-\frac{x^{2}}2+o(x^{2})$\\
$\ln(1+x)$ & $x-\frac{x^{2}}2+\frac{x^{3}}3+o(x^{3})$\\
$(1+x)^{\alpha}$ & $1+\alpha x+\frac{\alpha(\alpha-1)}{2}x^{2}+o(x^{2})$\\
$\tan x$ & $x+\frac{x^{3}}3+o(x^{3})$\\
$\arcsin x$ & $x+\frac{x^{3}}6+o(x^{3})$\\
$\arctan x$ & $x-\frac{x^{3}}3+o(x^{3})$\\
$\sqrt{1+x}$ & $1+\frac{x}{2}-\frac{x^{2}}8+o(x^{2})$\\
$\dfrac{1}{1-x}$ & $1+x+x^{2}+o(x^{2})$\\
\bottomrule
\end{tabularx}
\end{center}

\begin{center}
\footnotesize
\renewcommand{\arraystretch}{1.25}
\begin{tabularx}{\textwidth}{@{}>{\raggedright\arraybackslash}p{4.4cm} X@{}}
\toprule
\textbf{等价无穷小} & \textbf{公式}\\
\midrule
$\sin u\sim u$、$\tan u\sim u$、$\arcsin u\sim u$、$\arctan u\sim u$ & （$u\to0$）\\
$\ln(1+u)\sim u$ & （$u\to0$）\\
$1-\cos u\sim\frac{u^{2}}2$ & （$u\to0$）\\
$(1+u)^{\alpha}\sim1+\alpha u$ & （$\alpha\neq0$）\\
$u-\sin u\sim\frac{u^{3}}6$、$\tan u-u\sim\frac{u^{3}}3$ & （$u\to0$，做差时用）\\
\bottomrule
\end{tabularx}
\end{center}

\begin{warnbox}[洛必达的三种失效（备用，遇到卡壳先换工具）]
\begin{enumerate}[label=(\alph*)]
  \item \textbf{越算越繁}：$\frac{\ln(1+x)}{x}$ 型，一行等价就完事。
  \item \textbf{循环}：$\frac{\ee^{x}}{\ee^{x}}$、$\frac{\sin x}{x}\ (x\to+\infty)$。
  \item \textbf{求导后极限不存在，但原极限存在}：$x\to+\infty$ 时
        $\frac{x-\sin x}{x+\sin x}\to1$，而洛必达给出 $\frac{1-\cos x}{1+\cos x}$（振荡、无极限），
        于是误判"不存在"。\textbf{这个反例的 $x\to+\infty$ 不能写成 $x\to0$}：
        $\lim\limits_{x\to0}\frac{x-\sin x}{x+\sin x}=0$（sympy 复核），两边答案完全不同。
\end{enumerate}
\end{warnbox}

\end{document}
